\documentclass[a4paper,12pt]{article}

\usepackage[T1]{fontenc}
\usepackage[italian]{babel}
\usepackage{graphicx}
\usepackage{geometry}
\usepackage{tabularx}
\usepackage[table]{xcolor}
\usepackage[most]{tcolorbox}
\usepackage{fancyhdr}
\usepackage[hidelinks]{hyperref}
\usepackage{tikz}
\usepackage{pgf-pie}

\newcommand{\groupmail}{alittlebyte.swe@gmail.com}
\newcommand{\grouplogo}{../../.github/site-src/assets/logo_alittlebyte.png}
\newcommand{\groupsymbol}{../../.github/site-src/assets/solo_logo.png}

\geometry{
    left=2.5cm,
    right=2.5cm,
    top=2.5cm,
    bottom=2.5cm,
    headheight=331pt,
    includeheadfoot
}

\pagestyle{fancy}
\fancyhf{}

\renewcommand{\headrulewidth}{0pt}
\fancyhead{}

\renewcommand{\footrulewidth}{0.4pt}
\fancyfoot[L]{\includegraphics[height=0.6cm]{\groupsymbol}\hspace{0.45em}aLittleByte}
\fancyfoot[R]{Glossario}

\fancypagestyle{firstpage}{
    \fancyhf{}
    \renewcommand{\headrulewidth}{0pt}
    \renewcommand{\footrulewidth}{0.4pt}
    \fancyhead[C]{%
        \makebox[\textwidth][c]{%
            \begin{tabular}{c}
                \includegraphics[height=10cm]{\grouplogo}\\[1ex]
                \small\texttt{\groupmail}
            \end{tabular}%
        }
    }
    \fancyfoot[L]{\includegraphics[height=0.6cm]{\groupsymbol}\hspace{0.45em}aLittleByte}
    \fancyfoot[R]{Glossario}
}

\begin{document}

\thispagestyle{firstpage}

\vspace*{0.5cm}

\begin{center}
    {\LARGE \textbf{Glossario}}
\end{center}

\vspace{1.5cm}

\begin{center}
    \renewcommand{\arraystretch}{1.3}
    \begin{tcolorbox}[
        width=0.92\textwidth,
        colback=gray!6,
        colframe=gray!45,
        boxrule=0.5pt,
        arc=2mm,
        boxsep=0pt,
        left=8pt, right=8pt, top=8pt, bottom=8pt
    ]
        \begin{tabularx}{\linewidth}{>{\bfseries}p{3.5cm} X}
            Gruppo      & aLittleByte (gruppo 19) \\
            Versione & 2.1.0 \\
            Redattore   & V. Y. Kaneda Fernandes, Leonardo Soligo, Eleonora Bellet \\
            Verificatori & Angela Maule, Eleonora Bellet, Diego Belluz\\
            Destinatari & Prof. Tullio Vardanega, Prof. Riccardo Cardin \\
            Data        & 26/08/2026
        \end{tabularx}
    \end{tcolorbox}
\end{center}

\newpage
\fancyfoot[R]{Glossario\hspace{0.8em}}
\newgeometry{left=2.5cm,right=2.5cm,top=2.5cm,bottom=2.5cm,includefoot}

\begin{center}
    {\Large \textbf{Registro delle Modifiche}}
\end{center}

\vspace{0.35cm}

\renewcommand{\arraystretch}{1.35}
\rowcolors{2}{gray!8}{white}
\begin{center}
    {\footnotesize
    \begin{tabularx}{\textwidth}{|l|l|X|p{2cm}|p{2cm}|p{2cm}|}
        \hline
        \rowcolor{gray!20}
        \textbf{Versione} & \textbf{Data} & \textbf{Descrizione} & \textbf{Redattore} & \textbf{Verificatore} & \textbf{Approvatore} \\
        2.1.0 & 26/08/2026 & Aggiunti termini della Specifica Tecnica. & - & - & - \\
        2.0.0 & 23/06/2026 & Approvazione versione 2.0.0 & - & - & Alex Oloni \\
        1.0.1 & 22/06/2026 & Aggiunti termini dell'Analisi dei Requisiti. & V. Y. Kaneda Fernandes & Eleonora Bellet & - \\
        1.0.0 & 28/05/2026 & Approvazione versione 1.0.0 & - & - & Alex Oloni \\
        0.4.0 & 22/05/2026 & Aggiunti termini del Piano di Qualifica, Analisi dei Requisiti e Norme di Progetto. & Leonardo Soligo & Angela Maule & - \\
        0.3.3 & 09/05/2026 & Aggiunti termini dell'Analisi dei Requisiti, Norme di Progetto e Piano di Progetto & Eleonora Bellet & Diego Belluz & - \\
        0.3.2 & 01/05/2026 & Aggiunti termini dell'Analisi dei Requisiti & Eleonora Bellet & Diego Belluz & - \\
        0.3.1 & 27/04/2026 & Aggiunta di alcuni termini delle Norme di Progetto e Piano di Progetto & Eleonora Bellet & Diego Belluz & - \\
        0.2.0 & 22/04/2026 & Aggiunta di alcuni termini della Analisi dei Requisiti& Leonardo Soligo & Eleonora Bellet & - \\
        0.1.1 & 12/04/2026 & Aggiunta termini al documento & V. Y. Kaneda Fernandes & Angela Maule & - \\
        0.1.0 & 30/03/2026 & Prima stesura del documento & V. Y. Kaneda Fernandes & Angela Maule & - \\
        \hline 
    \end{tabularx}
    }
\end{center}

\newpage
\pagenumbering{arabic}
\setcounter{page}{1}
\fancyfoot[R]{Glossario\hspace{0.8em}\thepage}
\newpage

\tableofcontents

\newpage


\section{A}

\subsection{Adapter}
Componente che traduce l'interfaccia di una tecnologia esterna nell'interfaccia richiesta dal dominio applicativo.

\subsection{API (Application Programming Interface)}
Insieme di operazioni e regole che permettono a componenti software diversi di comunicare.

\subsection{Architettura Esagonale}
Stile architetturale che separa la logica di dominio dalle tecnologie esterne tramite porte e adapter.

\subsection{Agile}
Modello di sviluppo basato su sprint bisettimanali, adottato dal team.

\subsection{AI Analyst}
Nucleo elaborativo del modulo AI Co-Pilot che esegue autonomamente classificazione documentale, OCR ed estrazione dati.

\subsection{AI Assistant Generativo}
Modulo che permette la creazione autonoma di contenuti aziendali (titolo, testo, immagini) a partire da un prompt.

\subsection{AI Co-Pilot per Consulenti del Lavoro (CdL)}
Sistema avanzato in grado di riconoscere automaticamente la tipologia di documenti caricati, estrarre i destinatari tramite OCR ed Entity Resolution.

\subsection{AI Post Generator}
Modulo esterno (attore di sistema) che costituisce il nucleo di intelligenza artificiale per l'elaborazione di testi, titoli e immagini.

\subsection{Amazon Bedrock}
Servizio AWS che rende disponibili modelli di fondazione (FM) tramite API, utilizzato per l'integrazione di funzionalità AI.

\subsection{Amministratore}
Ruolo responsabile della configurazione dell'ambiente standalone, della gestione dei privilegi e della definizione delle policy globali.

\subsection{Analisi dei Requisiti (AdR)}
Documento che formalizza le funzionalità del sistema, inclusa la modellazione degli attori, la descrizione dei casi d'uso e la classificazione dei requisiti.

\subsection{Analisi dei Rischi}
Processo strutturato per identificare, classificare e valutare potenziali eventi critici.

\subsection{Analisi Dinamica}
Attività di verifica del codice eseguita durante la sua esecuzione.

\subsection{Analisi Statica}
Attività di verifica del codice eseguita senza la sua esecuzione.

\subsection{Analista}
Ruolo incaricato dell'analisi dei requisiti e del dominio.

\subsection{Analista Dati}
Figura che supervisiona l'efficienza del sistema, monitorando le performance dell'IA attraverso feedback e metriche.

\subsection{Angular}
Framework open-source per lo sviluppo web.

\subsection{AWS}
Fornitore di servizi cloud utilizzati dal sistema per AI, storage, code e orchestrazione.

\subsection{AWS S3}
Nome del servizio AWS di object storage utilizzato per i file del sistema.

\subsection{AWS SQS}
Nome del servizio AWS di code di messaggi utilizzato per i task asincroni.

\subsection{Approccio Incrementale}
Metodo di stesura e sviluppo in cui i documenti e il prodotto vengono costantemente aggiornati e affinati nel tempo.

\subsection{Attore}
Ruolo ricoperto da un utente o da un sistema esterno che interagisce con l'applicazione.

\subsection{Attore Primario}
L'utente umano o il sistema esterno che avvia attivamente un caso d'uso.

\subsection{Attore Secondario}
Entità o sistema terzo coinvolto passivamente nel processo.

\section{B}
\subsection{Bookmark}
Funzionalità che permette di contrassegnare e salvare i prompt preferiti.

\subsection{Bozza (Draft)}
Stato temporaneo di un documento salvato prima della validazione o invio definitivo.

\subsection{Branch}
Ramificazione del repository Git per lo sviluppo isolato.

\subsection{Backend}
Parte server-side del sistema che espone le API e implementa la logica applicativa.

\subsection{Bootstrap}
Avvio e configurazione iniziale dell'applicazione e dei suoi servizi.


\section{C}

\subsection{Callback}
Risposta inviata a un orchestratore per comunicare l'esito di un'attività asincrona.
\subsection{Capitolato}
Documento tecnico fornito dal proponente che definisce requisiti e obiettivi.

\subsection{Caso d'Uso (UC)}
Descrizione di un'interazione tra un attore e il sistema.

\subsection{Cedolini Massivi}
File contenente le buste paga di più dipendenti, da suddividere automaticamente.

\subsection{Ciclo Iterativo}
Suddivisione del lavoro in periodi ricorrenti (sprint).

\subsection{Code Review}
Revisione del codice effettuata tra pari.

\subsection{CloudFront}
Servizio AWS di distribuzione dei contenuti statici tramite rete CDN.

\subsection{Container}
Ambiente isolato che esegue un'applicazione e le sue dipendenze.

\subsection{CI/CD (Continuous Integration/Continuous Deployment)}
Pratiche di integrazione, verifica e rilascio automatico del software.

\subsection{CDN (Content Delivery Network)}
Rete distribuita che avvicina agli utenti la distribuzione di contenuti statici.


\subsection{Committente (Proponente)}
Azienda che ha commissionato il progetto (Eggon S.r.l.).

\subsection{Confidenza (Confidence Score)}
Percentuale che indica l'affidabilità del riconoscimento automatico.

\subsection{Consulente del Lavoro (CdL)}
Attore principale che utilizza l'AI Co-Pilot per gestire documenti.

\subsection{Consuntivo}
Rendiconto delle ore effettivamente impiegate e dei costi sostenuti.

\subsection{Copie Atomiche}
Singole parti risultanti dallo split di un documento massivo.


\section{D}
\subsection{Dashboard}
Interfaccia grafica per visualizzare metriche e statistiche.

\subsection{Discord}
Piattaforma di comunicazione digitale utilizzata dal team per il coordinamento rapido, canali testuali/vocali e riunioni interne.

\subsection{Detex}
Strumento per rimuovere sintassi LaTeX ed estrarre testo puro.

\subsection{Diagramma delle classi UML}
Rappresentazione della struttura statica del sistema.

\subsection{Diario di Bordo}
Documento per il tracciamento settimanale delle attività.

\subsection{Dispaccio}
Processo automatico di invio dei documenti ai destinatari.

\subsection{Docker}
Piattaforma per la costruzione e l'esecuzione di applicazioni containerizzate.

\subsection{Docker Compose}
Strumento per orchestrare più container tramite una configurazione comune.

\subsection{Dompdf}
Libreria PHP per trasformare contenuti HTML in documenti PDF.

\subsection{DLQ (Dead Letter Queue)}
Coda che raccoglie i messaggi non elaborati dopo il numero massimo di tentativi.

\subsection{Dependency Injection}
Tecnica con cui le dipendenze di una classe vengono fornite dall'esterno invece di essere create al suo interno.

\section{E}

\subsection{Entity Resolution}
Processo di riconoscimento e associazione dei destinatari ai documenti.

\subsection{Esecuzione Asincrona}
Modalità di esecuzione di operazioni lunghe in background.

\subsection{ETag}
Identificatore della versione di una risorsa HTTP usato per evitare trasferimenti o rigenerazioni inutili.

\subsection{Eloquent}
ORM di Laravel utilizzato per mappare modelli applicativi e tabelle del database.

\subsection{Entity}
Oggetto di dominio dotato di identità e comportamento propri.

\subsection{Event-Driven}
Approccio in cui le componenti comunicano reagendo alla pubblicazione di eventi.

\section{F}
\subsection{Figma}
Strumento per la progettazione di mockup interattivi.

\subsection{Filament}
Framework PHP per interfacce amministrative, parte dello stack tecnologico.

\subsection{Filtro dinamico}
Strumento che aggiorna istantaneamente i dati visualizzati.

\subsection{Facade}
Interfaccia unificata che semplifica l'accesso a più componenti di un sottosistema.

\subsection{FPDF e FPDI}
Librerie PHP utilizzate rispettivamente per generare PDF e importare pagine da PDF esistenti.

\section{G}
\subsection{GDPR}
Regolamento europeo sulla protezione dei dati personali.

\subsection{Git}
Software per il controllo di versione distribuito.

\subsection{GitHub}
Piattaforma per hosting e gestione collaborativa di repository.

\subsection{Glossario}
Documento che definisce i termini tecnici e ambigui del progetto.

\subsection{GHCR (GitHub Container Registry)}
Registro GitHub per la pubblicazione delle immagini Docker prodotte dal progetto.

\subsection{Gateway}
Componente che espone un punto di accesso verso un servizio o sistema esterno.

\subsection{Quality Gate}
Controllo automatico che deve essere superato prima di integrare o rilasciare una modifica.

\section{H}
\subsection{Human-in-the-loop}
Approccio con revisione umana per casi a bassa confidenza dell'AI.

\subsection{Heartbeat}
Segnale periodico che indica a un orchestratore che un'elaborazione è ancora attiva.

\section{I}
\subsection{Indice di Gulpease (MPC-IG)}
Indice di leggibilità per testi in italiano.

\subsection{Intelligenza Artificiale (AI)}
Tecnologie che permettono ai computer di svolgere attività che richiederebbero l'intelligenza umana, come l'apprendimento, il riconoscimento di pattern, la comprensione del linguaggio naturale e la generazione di contenuti.

\subsection{Issue}
Strumento per tracciare task su GitHub.

\subsection{IaC (Infrastructure as Code)}
Gestione dell'infrastruttura tramite file di configurazione versionati e riproducibili.

\subsection{Idempotenza}
Proprietà per cui ripetere un'operazione produce lo stesso effetto della sua prima esecuzione.

\subsection{JSON}
Formato testuale per lo scambio di dati strutturati tra frontend e backend.

\subsection{JSONB}
Tipo binario di PostgreSQL per memorizzare dati JSON semi-strutturati.

\section{J}
\subsection{JWT (JSON Web Token)}
Token firmato utilizzato per rappresentare l'identità e autorizzare le richieste.

\section{K}
\subsection{KLOC}
Migliaia di linee di codice.

\subsection{KMS (Key Management Service)}
Servizio AWS per la gestione delle chiavi crittografiche usate nella cifratura dei dati.

\subsection{KPI (Key Performance Indicator)}
Indicatori numerici per misurare le performance del sistema.

\section{L}

\subsection{LaTeX}
Linguaggio di composizione per documenti professionali.

\subsection{Laravel}
Framework PHP per lo sviluppo del backend.

\subsection{Latenza}
Tempo di risposta del sistema a una richiesta.


\subsection{LLM (Large Language Model)}
Modello di linguaggio di grandi dimensioni per la generazione di testo.

\subsection{Lotto Documentale}
Insieme di file caricati simultaneamente nel sistema.

\subsection{LocalStack}
Emulatore locale dei servizi AWS utilizzato per lo sviluppo e i test senza credenziali reali.

\subsection{Listener}
Componente che reagisce a un evento pubblicato dal sistema, ad esempio per audit o metriche.

\subsection{LLM Gateway}
Interfaccia applicativa verso un modello linguistico esterno.

\section{M}
\subsection{Marcatore di trasparenza}
Dicitura "Creato da AI Assistant" inserita nei contenuti esportati.

\subsection{Metadati}
Informazioni di contesto associate a un file.

\subsection{Milestone}
Traguardo significativo nel cronoprogramma (RTB, PB).

\subsection{MinIO}
Storage compatibile con API S3 per la gestione dei file.

\subsection{Mitigazione (dei rischi)}
Azioni per ridurre probabilità o impatto di un rischio.

\subsection{Mockup}
Prototipo grafico dell'interfaccia utente.

\subsection{Modello C4}
Modello per descrivere l'architettura del sistema.

\subsection{Modello di Sviluppo}
Approccio metodologico adottato (Agile/SCRUM).

\subsection{Moduli Esterni}
Sistemi terzi con cui il software scambia dati (AI Analyst, AI Post Generator).

\subsection{MVVM (Model-View-ViewModel)}
Pattern che separa dati, logica di presentazione e interfaccia utente.

\subsection{MVP (Minimum Viable Product)}
Versione con funzionalità minime necessarie.

\section{N}
\subsection{Nexum}
Piattaforma digitale di Eggon per comunicazione e gestione HR.

\subsection{Norme di Progetto}
Documento che definisce Way of Working, ruoli, strumenti e procedure.

\subsection{Nginx}
Web server e reverse proxy utilizzato per servire la SPA e inoltrare le richieste a PHP-FPM.

\subsection{Node.js}
Runtime JavaScript utilizzato dagli strumenti di build e sviluppo del frontend.

\section{O}
\subsection{OCR (Riconoscimento Ottico dei Caratteri)}
Tecnologia per estrarre testo da immagini o PDF scannerizzati.

\subsection{OpenAPI}
Specifica standard che descrive il contratto delle API REST e permette di generare client tipizzati.

\subsection{OpenTelemetry (OTel)}
Standard per raccogliere e trasportare metriche, tracce e log applicativi.

\subsection{Orval}
Strumento che genera automaticamente il client TypeScript a partire da una specifica OpenAPI.

\subsection{ORM (Object-Relational Mapping)}
Tecnica che rappresenta le tabelle di un database tramite oggetti del linguaggio di programmazione.

\subsection{OTLP (OpenTelemetry Protocol)}
Protocollo usato da OpenTelemetry per inviare segnali osservabili a un Collector.

\section{P}
\subsection{PB (Product Baseline)}
Fase finale del progetto per lo sviluppo e la consegna del prodotto.

\subsection{Penpot}
Strumento open source scelto per la progettazione di mockup.

\subsection{Pest}
Framework di testing per PHP.

\subsection{Pianificazione}
Attività di definizione di compiti, scadenze e responsabilità.

\subsection{Piano di Progetto}
Documento con pianificazione, rischi, preventivi e consuntivi.

\subsection{Piano di Qualifica}
Documento con metriche, test e obiettivi di qualità.


\subsection{Postcondizioni}
Stato finale del sistema dopo un caso d'uso.

\subsection{PostgreSQL}
Database relazionale del sistema.

\subsection{Porta}
Interfaccia astratta attraverso cui il dominio comunica con un adapter primario o secondario.

\subsection{PHP-FPM}
Gestore di processi PHP che esegue l'applicazione Laravel tramite FastCGI.

\subsection{PHP}
Linguaggio utilizzato per lo sviluppo del backend Laravel.

\subsection{SQL (Structured Query Language)}
Linguaggio utilizzato per definire e interrogare i dati del database relazionale.

\subsection{Prometheus}
Sistema per la raccolta e l'archiviazione di metriche esposte dai servizi.

\subsection{PSR-20}
Standard PHP che definisce l'interfaccia comune per un orologio astratto.

\subsection{Proxy Inverso}
Server che riceve le richieste dei client e le inoltra ai servizi interni appropriati.

\subsection{Precondizioni}
Stato obbligatorio prima di un caso d'uso.

\subsection{Preventivo}
Stima delle ore e dei costi prima dello sprint.

\subsection{Progettista}
Ruolo responsabile dell'architettura del sistema.

\subsection{Programmatore}
Ruolo responsabile della codifica.

\subsection{Prompt}
Testo fornito all'AI per generare un contenuto.

\subsection{Pull Request}
Richiesta di integrazione di modifiche su GitHub.

\subsection{PWA (Progressive Web App)}
Applicazione web con esperienza simile a un'app nativa.

\section{R}
\subsection{Rating}
Punteggio assegnato dall'utente al contenuto generato.

\subsection{RBAC (Role-Based Access Control)}
Controllo degli accessi basato sui ruoli.

\subsection{Redattore}
Utente che utilizza l'AI Assistant per creare comunicazioni.

\subsection{Redis}
Sistema in-memory utilizzato per cache, sessioni e rate limiting; non gestisce la coda della pipeline.

\subsection{Requisiti Funzionali}
Specifiche che descrivono i comportamenti, i servizi e le funzioni che il sistema deve eseguire per soddisfare i bisogni dell'utente.

\subsection{Requisiti opzionali}
Funzionalità accessorie che possono essere implementate solo se le risorse e i tempi lo permettono, senza impattare sugli obiettivi primari

\subsection{Responsabile}
Ruolo di project manager: coordina, pianifica e approva.

\subsection{Repository}
Archivio digitale per documenti e codice.

\subsection{Rischio}
 Evento o condizione incerta la cui occorrenza può influire (positivamente o negativamente) su costi, tempi o qualità degli obiettivi di progetto.

\subsection{RTB (Requirements and Technology Baseline)}
Milestone di progetto finalizzata alla validazione dei requisiti analizzati e al consolidamento dell'architettura tecnologica tramite prototipazione (PoC).

\subsection{REST (Representational State Transfer)}
Stile per la progettazione di servizi web basati su risorse e metodi HTTP.

\subsection{RxJS}
Libreria TypeScript per la programmazione reattiva tramite Observable e stream asincroni.

\subsection{Reactive Forms}
API Angular per definire e validare form tramite oggetti e controlli espliciti.

\section{S}
\subsection{Scenario Alternativo}
Percorso di gestione di errori o eccezioni.

\subsection{Scenario Principale}
Sequenza ottimale di interazioni.

\subsection{Soglia di Confidenza}
Valore limite per l'intervento manuale (es. 80\%).

\subsection{Sottocaso}
Scomposizione di un caso d'uso complesso.

\subsection{Split Automatico}
Suddivisione di documenti multi-destinatario.

\subsection{S3 (Simple Storage Service)}
Servizio AWS di object storage per documenti, sotto-documenti e asset statici.

\subsection{SQS (Simple Queue Service)}
Servizio AWS di code di messaggi usato per distribuire i task asincroni ai worker.

\subsection{SSE (Server-Sent Events)}
Tecnologia che consente al server di inviare aggiornamenti in tempo reale al browser tramite una connessione HTTP persistente.

\subsection{State Machine}
Modello che descrive stati e transizioni di una procedura orchestrata, come una pipeline asincrona.

\subsection{Step Functions}
Servizio AWS che orchestra le fasi delle pipeline asincrone tramite una state machine.

\subsection{Strategy}
Pattern che incapsula algoritmi intercambiabili e ne consente la selezione a runtime.

\subsection{SPA (Single Page Application)}
Applicazione web che aggiorna le viste senza ricaricare interamente la pagina.

\subsection{TLS (Transport Layer Security)}
Protocollo crittografico utilizzato per proteggere le comunicazioni HTTPS.

\subsection{Sprint}
Periodo di lavoro di breve durata (tipicamente 1-2 settimane) durante il quale il team realizza un insieme definito di attività e obiettivi

\subsection{Stakeholder}
Soggetti con interesse nel progetto (committente, docenti, utenti).

\subsection{Standalone}
Applicazione che funziona in modo indipendente.

\subsection{Stile Editoriale}
Parametro che calibra lo stile di scrittura dell'AI.

\section{T}
\subsection{Test di Accettazione}
Test che validano il sistema rispetto al capitolato.

\subsection{Test di Integrazione}
Test tra più moduli.

\subsection{Test di Regressione}
Test per rilevare nuovi errori dopo modifiche.

\subsection{Test di Sistema}
Test end-to-end in ambiente simulato.

\subsection{Test di Unità}
Test su singole funzioni in isolamento.

\subsection{Timeout}
Interruzione di un'operazione che supera il tempo massimo.

\subsection{Textract}
Servizio AWS di OCR che estrae testo e dati strutturati da documenti.

\subsection{Terraform}
Strumento per descrivere e creare l'infrastruttura tramite Infrastructure as Code.

\subsection{Tenant}
Organizzazione o spazio logico isolato i cui dati sono identificati da un tenant ID.

\subsection{TypeScript}
Superset tipizzato di JavaScript utilizzato per lo sviluppo del frontend Angular.

\subsection{Tono Comunicativo}
Parametro che calibra il tono dell'AI.

\subsection{Trigger}
Evento che innesca un caso d'uso.

\section{U}
\subsection{Utente Operativo}
Utente autorizzato alle funzionalità ordinarie.

\subsection{UUID (Universally Unique Identifier)}
Identificatore progettato per essere univoco, usato per nominare risorse e file.

\subsection{Observable}
Astrazione RxJS che rappresenta una sequenza di valori prodotti nel tempo.

\section{V}
\subsection{Verbale Esterno}
Resoconto di riunione con il proponente.

\subsection{Verbale Interno}
Resoconto di riunione del team.

\subsection{Verificatore}
Ruolo responsabile della verifica di documenti e codice.

\subsection{Versioning}
Sistema che tiene traccia delle versioni.

\subsection{Value Object}
Oggetto di dominio definito dal proprio valore, privo di identità autonoma.

\section{W}
\subsection{Walkthrough}
Metodo di revisione informale di documenti o codice.

\subsection{Way of Working (WoW)}
Insieme di norme, strumenti e procedure del gruppo.

\subsection{Workflow}
Sequenza automatica di passaggi per l'elaborazione.

\subsection{Worker}
Processo che consuma task da una coda ed esegue le operazioni asincrone assegnate.

\end{document}
